\scanpage{52}
\mspage{\handwritten{20}}
% archive sheet number: 51

\deleted{Notons $\handwritten{\varphi}_i:G\handwritten{\longrightarrow}\text{\uline{GL}}_k(V_i)$ le morphisme déterminé par l'opération de $G$ sur $V_i$, et $\handwritten{\varphi}:G\handwritten{\longrightarrow}\text{\uline{GL}}_k(A)$ le morphisme déterminé par l'opération de $G$ sur $A$. Alors, puisque $A$ est réunion des $V_i$, $\operatorname{Ker}\handwritten{\varphi}$ est int\deleted{re}\inserted{er}section des $\operatorname{Ker}\handwritten{\varphi}_i$; or $\operatorname{Ker}\handwritten{\varphi}$ est le $k$-groupe unité; puisque $G$ est n\oe{} thérien, il existe $i\in I$ tel que $\operatorname{Ker}\handwritten{\varphi}_{\unsure{i|1}}$ soit le $k$-groupe unité, ce qui implique (\deleted{\unsure{I|II}} 1.4.2) que $\handwritten{\varphi}_i$ est une immersion fermée, c.q.f.d.}

\uline{Lemme 11.12.- Soient $k$ un corps, $G$ un $k$-groupe quasi-compact. Posons $A=\handwritten{\mathscr{A}}(G)$. Etant donné $x\in A$, il existe \deleted{un sous-$k$-espace vectoriel de dimension finie}\inserted{une sous-$k$-algèbre de type fini.} $V$ de $A$ tel que $x\in V$, que $\handwritten{\Delta}(V)\subset V\handwritten{\otimes}_k V$ et $u(V)\subset V$, où $u$ désigne l'involution de $A$ correspondant à la symétrie de $G$}

\inserted{D'après (11.3), on peut remplacer $G$ par $G_{\mathrm{af}}$, donc on peut supposer $G$ affine.}

Soit $(\handwritten{\varphi}_i)$ une base de $A$. Alors $\handwritten{\Delta}(x)=\handwritten{\sum\varphi}_i\handwritten{\otimes}a_i$ et $\handwritten{\Delta}(\handwritten{\varphi}_i)=\handwritten{\sum\varphi}_j\handwritten{\otimes}b_{ji}$; l'axiome (HA1) de (\deleted{Exp.} I 4.2) montre que $\handwritten{\sum\varphi}_i\handwritten{\otimes\Delta}(a_i)=\handwritten{\sum\Delta}(\handwritten{\varphi}_i)\handwritten{\otimes}a_i$, d'où $\handwritten{\Delta}(a_i)=\handwritten{\sum}b_{ij}\handwritten{\otimes}a_j$. Soit alors $V$ la \deleted{sous-espace vectoriel}\inserted{sous-$k$-algèbre} de $A$ engendrée par les $b_{ij}$ et les $u(b_{ij})$. Il est clair que $V$ est une $k$-algèbre de type fini. L'axiome (HA1) montre aussi que $\handwritten{\sum\Delta}(\handwritten{\varphi}_j)\handwritten{\otimes}b_{ji}=\handwritten{\sum\varphi}_j\handwritten{\otimes\Delta}(b_{ji})$, et on en déduit que $\handwritten{\Delta}(b_{ij})=\handwritten{\sum}b_{ik}\handwritten{\otimes}b_{kj}$ et que $\handwritten{\Delta}(u(b_{ij}))=\handwritten{\sum}u(b_{kj})\handwritten{\otimes}u(b_{ik})$. Puisque $\handwritten{\Delta}$ est un homomorphisme d'algèbres, on en déduit que $\handwritten{\Delta}(V)\subset V\handwritten{\otimes}_k V$ et $u(V)\subset V$. Enfin l'axiome (HA2) de (\deleted{Exp} I 4.2) montre que $a_i=\handwritten{\sum\eta}(a_j)b_{ij}$ et que $x=\handwritten{\sum\eta}(\handwritten{\varphi}_i)a_i$, si bien que $x\in V$.

\uline{Proposition 11.13.- \deleted{Soit $G$ un $S$-groupe affine} Soient $k$ un corps et $G$ un $k$-groupe affine d'algèbre $A$. . Alors $G$ est limite projective d'un système filtrant croissant de $k$-groupes affines de type fini, dont les morphismes de transition sont fidèlement plats}.

Comme l'ensemble des sous-$k$-algèbres de type fini de $A$ stables par $\handwritten{\Delta}$ et $u$ est stable par engendrement, $A$ est limite inductive d'une famille f ltrante croissante $(B_i)_{i\in I}$ de sous-$k$-algèbres de type fini stables par $\handwritten{\Delta}$ et $u$. Alors, cahque algèbre $B_i$, munie de l'ho-\linebreak momorphisme $B_i\handwritten{\longrightarrow}B_i\handwritten{\otimes}_k B_i$ déduit de $\handwritten{\Delta}$ et de la restriction de $u$ à $B_i$ est munie d'une s\unsure{\deleted{rt}\inserted{tr}|\deleted{tr}\inserted{rt}}ucture d'hyperalgèbre associative augmentée involutive, donc (\deleted{Exp.} I 4.2) est la $\overset{\handwritten{/}}{k}$-algèbre d'un $k$-groupe affine $G_i$, de type fini sur $k$. Enfin, puisque $A=\underset{\handwritten{\longrightarrow}}{\lim}V_{i\!\handwritten{\rightharpoondown}}$,
$\unsure{G|C}$
